\documentclass[12pt,a4paper]{article}  
\usepackage[utf8]{inputenc}
\usepackage[T1]{fontenc}
\usepackage[spanish]{babel}
\usepackage{geometry}
\usepackage{hyperref}
\usepackage{longtable}
\usepackage{graphicx}
\usepackage{enumitem}  
\usepackage{lmodern}

\geometry{margin=2.5cm}
\setlist[itemize]{topsep=4pt, itemsep=2pt, leftmargin=1.5em}

\title{Documentación del Proyecto: Simulación de Cafetería con Hilos y JavaFX}
\author{Autor: Joel Novo Pampillón}
\date{\today}

\begin{document}

\maketitle

\section*{Introducción}

Este proyecto consiste en una simulación de una cafetería usando Java y JavaFX. La idea principal es representar cómo se gestionan los pedidos de múltiples clientes al mismo tiempo usando hilos (Threads), de manera que la preparación y entrega de cafés se haga de forma eficiente y organizada.

Se introducen dos roles principales en la cafetería: los \textbf{baristas}, que preparan los cafés, y los \textbf{camareros}, que los sirven a los clientes. Para organizar la entrega de los cafés, se utiliza un \textbf{buffer}, un espacio donde los cafés preparados esperan a ser servidos. Todo esto permite visualizar de forma dinámica el flujo de clientes, preparación y entrega de cafés en tiempo real.


\section{Arquitectura del Sistema}

\subsection{Modelo}

\subsubsection*{Clases y responsabilidades}

\begin{itemize}
    \item \textbf{Cliente}
    \begin{itemize}
        \item Representa a un cliente con nombre y tiempo máximo de espera.
        \item Implementa \texttt{Runnable} para correr en su propio hilo.
        \item Pide un café y espera a que se lo sirvan, o se va enfadado si tarda demasiado.
    \end{itemize}

    \item \textbf{Camarero}
    \begin{itemize}
        \item Extiende \texttt{Thread}.
        \item Toma cafés listos del buffer y los entrega a los clientes.
        \item Si no hay cafés listos, espera hasta que un barista prepare uno.
    \end{itemize}

    \item \textbf{Barista}
    \begin{itemize}
        \item Extiende \texttt{Thread}.
        \item Prepara cafés para los clientes en orden.
        \item Coloca los cafés en el buffer para que los camareros los sirvan.
        \item Si el buffer está lleno, espera hasta que haya espacio.
    \end{itemize}

    \item \textbf{Cafetería}
    \begin{itemize}
        \item Gestiona la cola de clientes que esperan su café.
        \item Mantiene la sincronización entre clientes, baristas y camareros.
        \item Controla si la cafetería está abierta o cerrada.
        \item Envía eventos al controlador para actualizar la interfaz.
    \end{itemize}

    \item \textbf{Buffer de cafés}
    \begin{itemize}
        \item Espacio donde se colocan los cafés ya preparados.
        \item Evita que los baristas sobreproduzcan y que los camareros sirvan cafés que no están listos.
        \item Tiene capacidad limitada para simular la “mesa” de cafés listos.
    \end{itemize}
\end{itemize}

\paragraph{Conceptos de programación concurrente:}
\begin{itemize}
    \item Uso de \texttt{synchronized}, \texttt{wait()} y \texttt{notifyAll()} para coordinar los hilos.
    \item Se asegura que los clientes solo reciban su café si está listo y dentro de su tiempo de espera.
    \item Evita que los baristas o camareros trabajen sobre recursos que no están disponibles.
\end{itemize}


\subsection{Vista (\texttt{cafeteria.fxml})}

La interfaz está desarrollada con JavaFX y definida mediante FXML:

\begin{itemize}
    \item Un área para mostrar los eventos en tiempo real.
    \item Contadores para clientes contentos y enfadados.
    \item Un botón para iniciar la simulación.
\end{itemize}

\paragraph{Ventajas:}
\begin{itemize}
    \item Separa la lógica del modelo de la visualización.
    \item Permite ver el flujo de clientes y cafés de manera clara.
\end{itemize}


\subsection{Controlador (\texttt{ControladorCafeteria.java})}

\begin{itemize}
    \item Inicializa la clase \texttt{Cafetería} y el buffer de cafés.
    \item Responde al botón de iniciar simulación.
    \item Crea y lanza los hilos de clientes, baristas y camareros.
    \item Actualiza la interfaz con los eventos usando \texttt{Platform.runLater}.
\end{itemize}

\paragraph{Flujo principal:}
\begin{enumerate}
    \item Se crean clientes y se agregan a la cola de la cafetería.
    \item Los baristas preparan los cafés y los colocan en el buffer.
    \item Los camareros toman los cafés del buffer y se los entregan a los clientes.
    \item Se registran eventos en la interfaz y se actualizan los contadores.
\end{enumerate}


\subsection{Launcher (\texttt{Launcher.java})}

\begin{itemize}
    \item Carga el archivo FXML desde los recursos.
    \item Muestra la ventana principal con título y tamaño configurados.
    \item Permite iniciar la simulación desde la interfaz.
\end{itemize}


\section{Casos de Uso}

\begin{longtable}{|c|l|l|l|}
\hline
\textbf{ID} & \textbf{Caso de Uso} & \textbf{Actor} & \textbf{Flujo Principal}  \\ \hline
CU1 & Entrar a la cafetería & Cliente & Llega, se agrega a la cola y espera café  \\ \hline
CU2 & Preparar café & Barista & Toma cliente de la cola, prepara café y lo coloca en el buffer \\ \hline
CU3 & Servir café & Camarero & Toma café del buffer y se lo entrega al cliente \\ \hline
CU4 & Recibir café & Cliente & Cliente recibe su café y se marca como contento \\ \hline
CU5 & Cliente se va & Cliente & Supera tiempo de espera y se retira enfadado \\ \hline
CU6 & Iniciar simulación & Usuario & Se crean hilos de clientes, baristas y camareros \\ \hline
CU7 & Cierre de cafetería & Sistema & Cafetería se cierra y todos los hilos finalizan \\ \hline
\end{longtable}


\section{Interfaz}

\subsection*{Elementos principales}

\begin{itemize}
    \item \textbf{Título de la ventana:} Cafetería Novotoast.
    \item \textbf{Área de eventos:} Muestra la llegada de clientes, preparación de cafés y entrega en tiempo real.
    \item \textbf{Contadores:} Clientes contentos y enfadados, actualizados automáticamente.
    \item \textbf{Botón de inicio:} Inicia la simulación y se desactiva mientras se ejecuta.
\end{itemize}

\subsection*{Diseño}

\begin{itemize}
    \item Disposición vertical: título, contadores, área de eventos y botón.
    \item Facilita seguir la simulación y ver cómo se mueven los clientes y los cafés.
\end{itemize}


\section{Conclusiones}

\begin{itemize}
    \item Se implementó una simulación de cafetería con hilos usando JavaFX.
    \item La introducción de baristas y buffer permite gestionar múltiples cafés de forma organizada.
    \item La coordinación entre clientes, baristas y camareros se realiza con \texttt{wait()} y \texttt{notifyAll()} para evitar conflictos.
    \item La interfaz muestra claramente la preparación y entrega de cafés, así como los clientes que se van contentos o enfadados.
\end{itemize}

\end{document}
